\documentclass{article}
\usepackage{amsmath, amssymb, amsthm}
\usepackage{hyperref}
\usepackage{geometry}
\geometry{margin=1in}

\title{Erd\H{o}s Problem \#961}
\date{Last updated: 2025-08-31}
\author{Source: erdosproblems.com}

\begin{document}
\maketitle

\noindent\textbf{Status:} OPEN \quad \textbf{No prize}

\noindent\textbf{Tags:} number theory

\medskip

\noindent\textbf{Problem Statement:}

Let $f(k)$ be the minimal $n$ such that every set of $n$ consecutive integers $&#62;k$ contains an integer divisible by a prime $&#62;k$. Estimate $f(k)$.




In other words, how large can a consecutive set of $k$-smooth integers be? Sylvester and Schur (see \cite{Er34}) proved $f(k)\leq k$ and Erd\H{o}s \cite{Er55d} proved\[f(k)&lt;3\frac{k}{\log k}.\]Jutila \cite{Ju74} and Ramachandra, and Shorey \cite{RaSh73} proved\[f(k) \ll \frac{\log\log\log k}{\log \log k}\frac{k}{\log k}.\]It is likely that $f(k) \ll (\log k)^{O(1)}$. This is essentially equivalent to [683] .




References

[Er34] Erd\H{o}s, Paul, A Theorem of Sylvester and Schur . J. London Math. Soc. (1934), 282--288.

[Er55d] Erd\H{o}s, P., On consecutive integers . Nieuw Arch. Wisk. (3) (1955), 124--128.

[Ju74] Jutila, Matti, On numbers with a large prime factor. {II} . J. Indian Math. Soc. (N.S.) (1974), 125--130.

[RaSh73] Ramachandra, K. and Shorey, T. N., On gaps between numbers with a large prime factor . Acta Arith. (1973), 99--111.

\noindent\textbf{Related OEIS sequences:} A213253


\bigskip
\noindent\small{Source: \url{https://www.erdosproblems.com/961}}
\end{document}
